\section{Background and Related Work}

\subsection{Machine Learning Operations (MLOps)}

MLOps is an emerging discipline that combines machine learning, software engineering, and DevOps practices to streamline the development, deployment, and maintenance of ML systems in production \cite{sculley2015hidden}. Key principles include:

\subsubsection{Reproducibility}
Every aspect of the ML pipeline—data, code, configurations, and environment—must be version-controlled to ensure consistent results across different executions and environments.

\subsubsection{Automation}
Manual processes introduce errors and inefficiency. Automated pipelines for training, testing, and deployment reduce human intervention and accelerate iteration cycles.

\subsubsection{Monitoring}
Production models require continuous monitoring for data drift, model degradation, and system performance to maintain reliability over time.

\subsubsection{Collaboration}
ML projects involve data scientists, ML engineers, and DevOps teams. Proper tooling and workflows facilitate collaboration and knowledge sharing.

\subsection{MLOps Tools and Frameworks}

\subsubsection{DVC (Data Version Control)}
DVC extends Git's version control capabilities to large datasets and models. It tracks data changes, enables reproducible pipelines, and integrates with remote storage \cite{dvc2020}.

\subsubsection{MLflow}
MLflow is an open-source platform for managing the ML lifecycle, including experiment tracking, reproducible runs, and model deployment \cite{zaharia2018mlflow}. Its key components are:
\begin{itemize}
    \item \textbf{Tracking}: Logs parameters, metrics, and artifacts
    \item \textbf{Projects}: Packages code in reusable format
    \item \textbf{Models}: Standardized format for packaging models
    \item \textbf{Registry}: Central repository for model versioning
\end{itemize}

\subsubsection{FastAPI}
FastAPI is a modern Python web framework optimized for building APIs with automatic validation, documentation, and high performance \cite{fastapi2021}. It leverages Python type hints and Pydantic for data validation.

\subsubsection{Docker and Containerization}
Containerization ensures consistent runtime environments across development, testing, and production. Docker packages applications with all dependencies, eliminating "works on my machine" problems \cite{merkel2014docker}.

\subsection{Geospatial Machine Learning}

Geospatial ML involves analyzing data with geographic coordinates. Common techniques include:

\begin{itemize}
    \item \textbf{Distance calculations}: Haversine formula for GPS coordinates
    \item \textbf{Trajectory analysis}: Speed, acceleration, bearing from GPS traces
    \item \textbf{Proximity features}: Distance to points of interest (bus stops, buses)
    \item \textbf{Spatial clustering}: Grouping geographic patterns (DBSCAN, HDBSCAN)
\end{itemize}

\subsection{Clustering Algorithms}

\subsubsection{DBSCAN}
Density-Based Spatial Clustering of Applications with Noise (DBSCAN) groups points based on density, identifying clusters of arbitrary shape and marking outliers as noise \cite{ester1996density}.

\subsubsection{HDBSCAN}
Hierarchical DBSCAN extends DBSCAN by constructing a cluster hierarchy and extracting optimal clusters at varying density levels \cite{campello2013density}. Advantages include:
\begin{itemize}
    \item No need to specify number of clusters
    \item Automatic noise detection
    \item Handles varying density clusters
    \item More robust parameter selection
\end{itemize}

\subsection{Related Work}

Transportation behavior classification using smartphone sensors has been explored in various contexts:

\begin{itemize}
    \item \textbf{Mode detection}: Identifying transportation modes (walking, bus, train) using accelerometer and GPS \cite{stenneth2011transportation}
    \item \textbf{Activity recognition}: Classifying user activities from sensor fusion \cite{shoaib2014survey}
    \item \textbf{Bus ridership}: Estimating passenger counts using Wi-Fi and Bluetooth signals \cite{kostakos2009brief}
\end{itemize}

However, few works focus on the complete MLOps lifecycle for such systems, making this project a comprehensive case study in production ML engineering.
